\section{Experimental protocols}
\label{sec:experimental-design}

We distinguish completed analyses, reported in Section~\ref{sec:results}, from
additional experiments required to evaluate the statistical properties of the
method. The designs below specify the data-generating mechanisms, sample sizes,
comparators, estimands, replication, and primary analyses \emph{before} results
are generated. Every figure and table labeled \textbf{Planned} is an unfilled
result location, not an observation or a prediction of performance. The
complete configuration, data-source, metric, and reporting specifications are
maintained with the implementation in the accompanying reproducibility
protocol; the scientific definitions and primary comparisons are given here.

\subsection{Common evaluation conventions}
\label{subsec:evaluation-conventions}
All simulations use independent datasets as the unit of replication. A
configuration is fixed before the generating seed is drawn; all methods in a
configuration receive identical simulated observations and candidate boundaries.
The default generation scheme uses NumPy \texttt{SeedSequence} with separate
components for experiment, design cell, and replicate, so new cells do not
change existing seeds. The primary comparison uses 100 independently generated
datasets per cell unless otherwise specified below. A five-dataset pilot per
cell checks feasibility and memory use; pilot results are not included in
reported performance estimates. Tests of the exact recursions use separate
non-random exhaustive enumeration and numerical-integration oracles.

For boundary evaluation, use a maximum-cardinality one-to-one matching of
estimated to generating boundaries within tolerance $\tau$. Among matchings
with maximal cardinality, minimize total displacement. The primary tolerance
is $\tau=3$ indices; report $\tau\in\{0,1,5\}$ as sensitivity analyses.
Compute precision, recall, and $F_1$ from the matched pairs, counting unmatched
estimated and true boundaries in the usual manner. Matched-location error is
not defined when no boundaries match. The segment-count metrics are
$\Pr(\widehat k=k^\star)$, $|\widehat k-k^\star|$, and the probability of
$k=k^\star$ under the count posterior. For signal estimation report mean
squared error relative to the \emph{generating segment functional}, not across
families with incompatible response scales.

For posterior probabilities $\widehat p_{ri}=p(b_i=1\mid y^{(r)})$, report
Brier score, logarithmic score using a prespecified numerical clipping rule,
and reliability curves with counts per bin. Use ten equal-mass bins as the
primary diagnostic, separately reporting event prevalence; show equal-width
bins as a sensitivity analysis. Confidence intervals for simulation metrics
use the independent sequence as the resampling unit (2,000 stratified
bootstrap replicates with the resampling seed declared in the protocol).
Dependent positions from one sequence are not treated as independent
replicates. Intervals are descriptive across generated datasets, not
frequentist coverage statements about an individual posterior interval.

For any algorithm fitted to held-out data, hyperparameters and comparator
penalties are selected on disjoint training/validation sequences or by a
predeclared training-only rule. Methods are compared at matched input,
coordinate axis, and observation law. An agreement score against the
\BayesBreak MAP is never labeled ground-truth accuracy. Unless the models
define proper probabilities for the same observed variables and base measure,
their marginal-likelihood values are not directly compared.

\begin{table}[!htbp]
\centering\small
\caption{Planned validation matrix. Entries specify experiments to run; they are
not additional results. Exact parameter grids and artifact destinations are
included in the companion experimental protocol.}
\label{tab:planned-experiment-matrix}
\begin{tabularx}{0.98\textwidth}{@{}L{0.07\textwidth}L{0.28\textwidth}Y L{0.21\textwidth}@{}}
\toprule
ID & Question & Main design and primary estimand & Intended output \\
\midrule
E00 & Exactness of partition inference & Exhaustive partitions for $n\le10$ and integrated-block quadrature & Oracle agreement table \\
E01 & Posterior calibration and recovery & Gaussian $n\in\{256,512\}$, $k^\star\in\{1,3,6\}$, varying jump size & Reliability curves and recovery table \\
E02 & Observation-model generality & Gaussian, Poisson, Binomial, negative Binomial, and Beta-valued responses & Family recovery table \\
E03 & Irregular-design prior & Uniform, gap-hazard, length-cohesion, combined priors; controlled gaps & Boundary mass and detection plots \\
E04 & Common-changepoint pooling & $S\in\{1,4,16,64\}$ sequences, with and without shared truth & Recovery and calibration table \\
E05 & Latent-template estimation & $G\in\{1,2,3\}$; separation/noise sweeps; multiple restarts & ARI and collapse plots \\
E06 & Nonconjugate segment integration & Logistic-normal blocks, adaptive quadrature reference & Error versus block length, posterior discrepancy \\
E07 & Misspecification and missingness & Heavy tails, overdispersion, absent outcomes & Coverage and robustness plots \\
E08 & Posterior prediction & Disjoint held-out contiguous regions; correct family scoring & Test log score and uncertainty table \\
E09 & External comparators and case studies & Matched coordinate/likelihood analyses and annotations when verified & Independent benchmark tables \\
E10 & Scaling and resource cost & $n\in\{128,256,512,1024,2048\}$ and $k_{\max}\in\{5,10,20\}$ & Time and memory plots \\
\bottomrule
\end{tabularx}
\end{table}

\subsection{E00: exactness checks}
For $n\in\{2,\ldots,10\}$, enumerate every admissible partition at every
segment count. Test uniform and nonuniform normalized priors, including zero
prior factors and minimum-segment constraints. Compare $L_{k,n}$, $R_{k,0}$,
$C_k$, $p(k\mid y)$, ordered and event boundary marginals, posterior moments,
and conditional and global MAP decisions against direct enumeration. Include a
Gaussian example with negative segment posterior means to test signed moment
aggregation. Independently compare the Gaussian, Poisson, Binomial, and
negative-Binomial block integrals to one-dimensional numerical integration for
unit and nonunit exposure/trial/precision descriptors. Structural infeasibility
is not conflated with a block whose observed-data likelihood equals one.

\begin{table}[!htbp]
\centering\small
\caption{\textbf{Table P0 (E00).} Independent verification of the exact
recursions and family-specific integrals (artifact \texttt{03a4d1c38bfe08c1}).
Maximum observed numerical discrepancy across 380 enumerated test cases; all 380 passed.}
\label{tab:planned-oracles-paper}
\begin{tabularx}{0.88\textwidth}{@{}Y L{0.17\textwidth}L{0.17\textwidth}L{0.19\textwidth}@{}}
\toprule
Quantity & Cases tested & Maximum absolute error & Acceptance tolerance \\
\midrule
Partition sums and normalizers & 4 families; $n\in\{2,\ldots,10\}$; nonuniform prior factors (Gaussian) & $1.07\times10^{-14}$ & $10^{-8}$ on log scale \\
Boundary probabilities and moments & 4 families; $n\in\{2,\ldots,10\}$; first moments; 2nd-moment vs.\ quadrature (Gaussian, Poisson) & $7.11\times10^{-15}$ prob.; $6.30\times10^{-7}$ 2nd-moment & $10^{-9}$ prob.; $10^{-6}$ 2nd-moment \\
Family-specific integrated likelihoods & 4 families; $n\in\{2,\ldots,10\}$; unit and nonunit descriptors vs.\ 1D adaptive quadrature & $3.91\times10^{-14}$ & $10^{-6}$ (absolute) \\
\bottomrule
\end{tabularx}
\end{table}

\subsection{E01 and E02: simulated recovery, calibration, and family generality}
E01 uses regularly spaced Gaussian sequences of lengths $256$ and $512$, with
$k^\star\in\{1,3,6\}$, minimum generating segment length 32, and adjacent
mean differences $\Delta/\sigma\in\{0.5,1,2\}$. For $k^\star=1$, the
jump-size parameter is immaterial and duplicate null cells are omitted.
Use the fixed proper segment prior $\mu\sim N(0,4)$ and known
observation variance $\sigma^2=1$; these quantities are held constant
for the confirmatory runs. Analyze both the
model-averaged boundary-event posterior and the posterior conditional on the
true segment count as separately named outputs. Primary metrics are boundary
$F_1@3$, count posterior coverage, Brier score, and calibration by boundary
probability. Do not pool boundary positions as independent bootstrap units.

E02 fixes three generating segments and $n=240$ for conjugate families; an
initial $n=120$ setting is used for Beta-valued quadrature because it requires
numerical integration for every candidate segment. The exact generating
parameters are specified as rates, probabilities, dispersion and observation
precisions, with two effect-size settings per family. Report each family on
its own natural signal scale and include a standardized estimation metric for
display across panels. Use 100 simulated datasets per conjugate cell
and 50 per quadrature cell, after feasibility checks. The comparison tests
the common inference algorithm, not whether one observation family is superior
to another.

\begin{figure}[!htbp]
\centering
\BBPlannedPanel{(a) Boundary calibration}{Empirical frequency}{Predicted probability}
\hfill
\BBPlannedPanel{(b) Recovery versus signal size}{$F_1@3$}{$\Delta/\sigma$}
\caption{\textbf{Planned Figure P1 (E01).} Posterior calibration and detection
versus signal-to-noise ratio. Show independent-dataset uncertainty bands and
per-bin sample counts in the final figure. No numerical result is asserted here.}
\label{fig:planned-calibration-paper}
\end{figure}

\begin{table}[!htbp]
\centering\small
\caption{\textbf{Planned Table P1 (E01--E02).} Predeclared reporting columns.
One row per experiment, design cell, and method will be generated by the
simulation registry. The values below have not been measured.}
\label{tab:planned-recovery-paper}
\begin{tabularx}{0.94\textwidth}{@{}L{0.29\textwidth}L{0.15\textwidth}L{0.15\textwidth}L{0.15\textwidth}Y@{}}
\toprule
Design and method & $N$ datasets & $F_1@3$ [95\% CI] & $\Pr(\widehat k=k^\star)$ & Brier / log score \\
\midrule
Gaussian, varying SNR & TBD & TBD & TBD & TBD \\
Poisson / Binomial / NegBin & TBD & TBD & TBD & TBD \\
Beta-valued quadrature & TBD & TBD & TBD & TBD \\
\bottomrule
\end{tabularx}
\end{table}

\subsection{E03: irregularly spaced observations}
Simulate $n=240$ design points with three piecewise-constant segments and an
inserted empty interval between two clusters. Vary the relative gap width over
$\{1,4,16\}$ times the within-cluster spacing and the changepoint location
relative to the large gap. Analyze the same data under (i) index-uniform
partitions, (ii) a candidate-boundary hazard proportional to the physical gap,
(iii) segment-length cohesion, and (iv) their product. Normalize each
conditional partition prior with its own $C_k$ and report posterior boundary
mass at and near the gap as well as one-to-one recovery. Report a separate
simulation in which the generating boundary is \emph{not} at the large gap so
that the prior's false-positive behavior can be measured. This comparison is
a prior-sensitivity analysis, not a guarantee that any gap prior is uniformly
preferable.

\begin{figure}[!htbp]
\centering
\BBPlannedPanel{(a) Change at the long gap}{Posterior boundary mass}{Relative gap width}
\hfill
\BBPlannedPanel{(b) Change away from the gap}{False-positive probability}{Relative gap width}
\caption{\textbf{Planned Figure P6 (E03).} Comparison of index-uniform,
candidate-gap, length-cohesion, and combined priors under both generating
boundary placements. Show uncertainty across independent simulated datasets.}
\label{fig:planned-irregular-paper}
\end{figure}

\subsection{E04 and E05: hierarchical segmentation}
For E04, generate $S\in\{1,4,16,64\}$ independent Gaussian sequences
conditionally on a common three-segment template, then replace 25\% and 50\%
of the sequences with shifted subject-specific boundaries in separate
misspecification cells. Compare the common-boundary posterior, separate fits,
and a suitable partial-sharing baseline on boundary $F_1@3$, calibration, and
signal error. Sequence-level bootstrap resampling respects the shared
generating configuration.

For E05, use $G\in\{1,2,3\}$ templates and $S\in\{24,48\}$ sequences, with
predeclared low, medium, and high separation cells. Run five initializations
per dataset and select using the \emph{recomputed final} finite-template
objective. Report label-invariant clustering accuracy and adjusted Rand index,
complete objective traces, the fraction of empty or duplicate templates, and
failures under reduced separation. The criterion and its MM algorithm are
those of Section~\ref{sec:extensions}; do not replace the criterion with a
normalized generative mixture during this experiment.

\begin{figure}[!htbp]
\centering
\BBPlannedPanel{(a) Pooled boundary recovery}{$F_1@3$}{Number of sequences}
\hfill
\BBPlannedPanel{(b) Template recovery}{Adjusted Rand index}{Separation level}
\caption{\textbf{Planned Figure P2 (E04--E05).} Effects of replication,
assumption violations, and between-template separation. Plot dataset-level
uncertainty and template-collapse frequency.}
\label{fig:planned-hierarchical-paper}
\end{figure}

\subsection{E06: nonconjugate numerical evidence}
Use intercept-only Binomial--logistic segments with normal prior on the logit,
block sizes $\{4,8,16,32,64,128\}$, and generating success probabilities
$\{0.05,0.5,0.95\}$. For each cell generate 100 independent blocks.
Establish a converged adaptive quadrature reference with absolute and relative
tolerances $10^{-11}$ and $10^{-10}$, checking agreement with a second numerical
method on a stress subset. Compare the segment log marginal likelihood and its
first two posterior signal moments. Include Laplace, the implemented
Jaakkola--Jordan bound, and expectation propagation only if its complete
normalizing constant and stopping rule pass validation. The current
\texttt{pg\_vb} option calls the JJ routine and is not counted as an
independent method. Propagate these same block-score arrays through a
segmentation simulation, reporting posterior count discrepancy, boundary
probability error, ranking changes, and runtime. The maximum over reachable
segments must be used when interpreting the uniform-error theorem; a
posterior-weighted error is a distinct diagnostic.

Results for this experiment are reported in Section~\ref{sec:results} (Figure~\ref{fig:planned-nonconjugate-paper}, Table~\ref{tab:planned-nonconjugate-paper}), artifact \texttt{9739f69c790aae6f}.

\subsection{E07 and E08: robustness and held-out prediction}
E07 repeats the $n=256$, three-segment recovery experiment with Student
$t_3$ noise fitted by a Gaussian block model, negative-Binomial counts fitted
by Poisson blocks, and overdispersed success counts fitted by Binomial blocks.
Separately remove 10\% or 30\% of responses completely at random, or a single
contiguous 15\% region. Missing outcomes are represented by the observation
mask, not a zero Gaussian precision or a zero Poisson exposure. Compare
boundary calibration, count selection, and frequentist recovery to the
correctly specified reference.

E08 holds out contiguous regions from synthetic and real sequences while
retaining the fitted coordinate endpoints. The split is made before any
hyperparameter estimation. Each fitted observation family uses its own
posterior-predictive distribution with the \emph{observed} new-data exposure,
trial count, or precision. Report log predictive density per held-out
observation and proper prediction intervals or randomized PIT diagnostics as
appropriate for the family. Distinguish prediction conditional on the
estimated partition from averaging over the segmentation posterior. A
pointwise composite score is never called a joint predictive density.

\begin{table}[!htbp]
\centering\small
\caption{\textbf{Planned Table P3 (E07--E08).} Robustness and predictive
performance. Comparisons will be paired within identical held-out splits.}
\label{tab:planned-robustness-paper}
\begin{tabularx}{0.94\textwidth}{@{}L{0.34\textwidth}L{0.22\textwidth}L{0.20\textwidth}Y@{}}
\toprule
Perturbation or model & $F_1@3$ [95\% CI] & Held-out log density / obs. & Prediction coverage \\
\midrule
Correctly specified reference & TBD & TBD & TBD \\
Heavy tails / overdispersion & TBD & TBD & TBD \\
Missingness or blocked holdout & TBD & TBD & TBD \\
\bottomrule
\end{tabularx}
\end{table}

\subsection{E09: external comparisons and real-data sensitivity}
Use the original well-log record, the probe-aligned array-CGH matrix, the
2015--2023 daily S\&P 500 series, and the \texttt{methylKit} chromosome-21
example with dataset identity and checksums fixed before fitting. Reproduce
the reported preprocessing and include the full well-log record alongside the
subsample. Array-CGH analysis uses per-probe precision weights derived from rolling median absolute deviation of log$_2$ ratio differences, with a pooled reference variance estimate across subjects; this is a data-driven specification rather than a physically calibrated noise model. The precision weights enter the Gaussian likelihood as known inverse-variance scale factors, not as generalized-Bayes likelihood powers. Sensitivity of the segmentation to the precision specification, and comparison with a uniform-precision analysis, remain as potential future analyses.
For methylation, compare count-based Binomial or Beta-Binomial modeling when
raw methylated and total reads are available to the existing continuous-Beta
response model; do not regard a Gaussian predictive fallback as a valid Beta
score. Vary $k_{\max}$ until posterior mass near the upper truncation is
small, or explicitly report truncation sensitivity if that criterion is not
met. Evaluate a grid of segment-count priors and physical-coordinate priors
with the normalizer recomputed for every model.

Compare with PELT, optimal partitioning, binary segmentation, wild binary
segmentation, and a Bayesian reference implementation when each comparator is
appropriate to the response and feasible at the target length. Distinguish
(a) reference-annotation accuracy when independent labels and alignment are
actually verified, (b) model-to-model agreement when only fitted boundaries
are available, and (c) held-out predictive comparisons on the same outcomes.
No figure caption or table may relabel model agreement as biological or
geological ground truth. Baseline penalties, feature transforms, and stopping
criteria are determined by disjoint tuning data or stated a priori.

\begin{figure}[!htbp]
\centering
\BBPlannedPanel{(a) Boundary agreement or accuracy}{$F_1@3$}{Method and dataset}
\hfill
\BBPlannedPanel{(b) Segment-count sensitivity}{$p(k\mid y)$}{Admissible $k$}
\caption{\textbf{Planned Figure P4 (E09).} Fair matched-axis comparisons and
posterior sensitivity to $k_{\max}$, distinguished by whether an independent
boundary reference is available.}
\label{fig:planned-external-paper}
\end{figure}

\begin{table}[!htbp]
\centering\small
\caption{\textbf{Planned Table P4 (E09).} Corrected real-data and comparator
study; every row will identify its observed variables, reference, and priors.}
\label{tab:planned-external-paper}
\begin{tabularx}{0.94\textwidth}{@{}L{0.23\textwidth}L{0.17\textwidth}L{0.18\textwidth}L{0.19\textwidth}Y@{}}
\toprule
Dataset / model & $\widehat k$, tail mass & Reference type & $F_1@3$ [95\% CI] & Test log density \\
\midrule
Well-log, full and subsampled & TBD & TBD & TBD & TBD \\
Array-CGH, corrected precision & TBD & TBD & TBD & TBD \\
S\&P 500 and methylation & TBD & TBD & TBD & TBD \\
\bottomrule
\end{tabularx}
\end{table}

\subsection{E10: time and memory complexity}
Benchmark $n\in\{128,256,512,1024,2048\}$ and
$k_{\max}\in\{5,10,20\}$ on a single fixed machine, using three warmup
runs and seven measured runs per configuration. Pin Python, NumPy, BLAS
thread counts and data types. Record construction of the segment-score array,
forward and backward recursions, MAP backtracking, posterior marginal
extraction, peak resident memory, and the total wall time separately. Include
the ordinary $\Theta(n^2)$ storage path and a documented streaming variant if
it is implemented. Report observed finite-range slopes and theoretical
worst-case complexity as different quantities.

Results for this experiment are reported in Section~\ref{sec:e10} (Figure~\ref{fig:E10-scaling-paper}, Table~\ref{tab:E10-scaling-paper}), artifact \texttt{bece422c9d90093f}.

\subsection{Criteria for replacing planned outputs}
A planned slot is populated only when the generating configuration, data or
simulation seed, code revision, software dependencies, method configuration,
metric definition, and output files are available together. A table that has
only some completed rows remains labeled planned until all reported entries
can be traced to executed measurements. The statistical results and their
limitations must be revised from those measurements, not from expected
outcomes. The retained findings in Section~\ref{sec:results} are historical
experiments with stated conditions; they are not the unfilled cells in the
planned validation matrix.
